%%%%

%%%   Самарский государственный технический университет

%%%   Кафедра прикладной математики и информатики

%%%

%%%   Преамбула для статьи в журнал "Вестник Самарского государственного

%%%   технического университета. Серия: «Физико-математические науки»"

%%%   Ориентирована под MikTEx 2.4 for Win32

%%%

%%%   Хотя сам журнал собирается под GNU/Linux на дистрибутиве TeXLive



\documentclass[11pt,a4paper,twoside]{article}



\usepackage[cp1251]{inputenc}

\usepackage[T2A]{fontenc}

\usepackage[english,russian]{babel}

\usepackage{samgtubib}

\usepackage[dvips]{graphicx}

\usepackage[multiple]{footmisc}



\usepackage{samgtu}

\renewcommand{\VestnikYear}{2009} % Год издания Вестника

\renewcommand{\VestnikNumber}{2\,(19)} % Номер Вестника

\renewcommand{\VestnikMonth}{Ноябрь} % Месяц выхода Вестника

\renewcommand{\VestnikData}{01.11.09} % Дата подписания Вестника в печать

\renewcommand{\VestnikUPL}{26,64} % Усл. печ. л.

\renewcommand{\VestnikUIL}{25,73} % Уч.-изд. л.

\renewcommand{\VestnikRegNumber}{479/09} % Регистрационный номер

\renewcommand{\VestnikOrder}{994} % Номер заказа







%

%

%

%\begin{document}



\renewcommand{\firstpage}{115}

\renewcommand{\lastpage}{123}





\udk{517.588:(517.958:536.212)} \msc{35Q80; 80A20, 33C70}





\title{Представление температурного поля для~полубесконечного тела, нагреваемого неподвижным лазерным лучом, через~гипергеометрические функции}% Полный заголовок

{Представление температурного поля для полубесконечного тела \dots}% Заголовок, который идёт в верхний колонтитул



\author{В.\,В.~Манако}{М\,а\,н\,а\,к\,о~В.\,В.}



\institute{\samgtu.}



\email{E-mail}{viktor.manako@mail.ru}



\otherinf{\fullauthor{Виктор Владимирович Манако} \regalia{к.ф.-м.н., доц.} \position{доцент} \dept{каф. общей физики и физики нефтегазового производства}}



\keywords{краевая задачи теплопроводности, круговой источник тепла, гипергеометрические ряды}



\workabstract{Рассматривается аналитическое выражение для нестационарного температурного поля в полубесконечном теле, нагреваемом круговым источником тепла, расположенным на свободной поверхности.

Нестационарное поле температур выражается через гипергеометрические  функции Аппеля  и~Сриваставы. 

Рассматриваются частные случаи нестационарного температурного поля, которые выражаются через функцию Кампе де Ферье.

Полученные выражения позволяют сразу отделить стационарную часть  температурного поля от нестационарной.

Рассчитаны стационарные температурные поля, образуемые круговым и гауссовым источниками. 

Существенных количественных различий в~этих полях не обнаружено. 

}



\engtitle{A representation in terms of hypergeometric functions for the temperature field in a semi-infinite body that is heated by a motionless laser beam}



\engauthor{V.\,V.~Manako}{M\,a\,n\,a\,k\,o~V.\,V.}



\otherenginfm{\fullauthor{Victor V. Manako}  \regalia{Ph. D. (Phys. \& Math.)} \position{Associate Professor} \dept{Dept. of General Physics and Physics of Oil and Gas Production}

}



\enginstitute{\engsamgtu.}



\engabstract{We have considered an analytical expression for the temperature field of a semi-infinite body that is heated by a circular heat source located at the free surface. Unsteady temperature field is expressed in terms of the Appell and the Srivastava hypergeometric functions. We have studied some special areas in heated body where a non-stationary temperature field is expressed in terms of the Kamp\'e de F\'eriet function. The obtained expressions have allowed to carry out the separation of the stationary and non-stationary parts of temperature field from each other. Calculations of the steady temperature fields generated by circular or Gaussian sources have been accomplished. Significant quantitative differences in these fields were not found.

}



\engkeywords{heat conductivity boundary value problem, circular heat source, hypergeometric series}



\date{11/XI/2011}

\revisiondate{16/IV/2012}





\maketitle 



\Section[n]{Постановка задачи} 

Пусть на поверхность $z=0$ однородного и изотропного тела, занимающего полупространство $z\geq0$, нормально падает поток лазерного излучения, создающий плоское пятно нагрева в~виде неподвижного круга радиусом ~$R$. 

Через это круглое пятно рассматриваемое тело получает подводимую извне энергию лазерного излучения, а~на остальной части поверхности $z=0$ теплообмен с~окружающей средой отсутствует (адиабатическая граница). 

Воспользуемся цилиндрической системой координат $(r,\,\varphi,\,z)$, у которой начало отсчёта совпадает с~центром пятна нагрева, ось $z$ совпадает с~осью лазерного луча и направлена в~глубь полубесконечного тела, а радиальная ось $r$ перпендикулярна лучу. Запишем уравнение теплопроводности

\begin{equation}

\label{monako:conductivity}

\frac{\partial^2T}{\partial r^2}+\frac{1}{r}\frac{\partial T}{\partial r}+\frac{\partial^2T}{\partial z^2}=\frac{1}{a}\frac{\partial T}{\partial t}

\end{equation}

и дополним его начальными и краевыми условиями:

\begin{equation}

\label{monako:conditions}

\frac{\partial T}{\partial z}\Big|_{z \to  0}=-\frac{q}{\lambda}\,H(R-r), \qquad 

\frac{\partial T}{\partial r}\Big|_{r \to  0}=T\big|_{r\to \infty}= T\big|_{z\to \infty}=T\big|_{t\to 0}=0,

\end{equation}

где $\lambda$ "--- теплопроводность, $a$ "--- температуропроводность, $H$ "--- единичная функция Хевисайда. 

Поверхностная плотность поглощённой мощности   выражается через коэффициент поглощения $A$ и полную мощность лазерного луча~$P$: $q=AP/(\pi R^2)$. 

Первое из условий \eqref{monako:conditions} подразумевает, что величина $ q $ внутри пятна нагрева постоянна и~равна нулю вне этого пятна. 



%Краевая задача \eqref{monako:conductivity}, \eqref{monako:conditions} представляет значительный интерес \cite{manako:pinsker}. 

В~настоящей работе решение краевой задачи \eqref{monako:conductivity}, \eqref{monako:conditions} представляется в~терминах специальных функций, при этом полагается, что величины  $ \lambda $, $ a $, $ q $ "--- постоянны.



\smallskip



\Section[n]{Дополнительные сведения} Обобщённой гипергеометрической функцией аргумента $ x $ называется ряд

\begin{equation}

\label{monako:hypergeom}

_{p}F_{q}\bigg[\,^{\displaystyle (a_p)}_{\,\displaystyle (b_q)}\,\,\bigg.\bigg|\,x\,\bigg]=\sum_{k=0}^\infty\frac{(a_1)_k(a_2)_k\dots(a_p)_k}{(b_1)_k(b_2)_k\dots(b_q)_k}\frac{x^k}{k!}=\sum_{k=0}^\infty \frac{\prod\limits_{\scriptscriptstyle j=1}^{\scriptscriptstyle p} (a_j)_k}{\prod\limits_{\scriptscriptstyle j=1}^{\scriptscriptstyle q} (b_j)_k}\frac{x^k}{k!}.

\end{equation}

Здесь $(a_j)_k =a_j(a_j+1) \dots (a_j+k-1)= \Gamma (a_j+k) / \Gamma (a_j)$ "--- символ Похгаммера; $\Gamma (x)$ "--- гамма-функция Эйлера. 

Множества чисел $(a_p)$ и $(b_q)$, содержащие соответственно $ p $ и $ q $ элементов, называются параметрами функции $_{p}F_{q}$.



Большинство стандартных специальных функций "--- частные случаи функции \eqref{monako:hypergeom}: 

\begin{equation}

\label{monako:BesselJ}

J_\nu (x)=\frac{1}{\Gamma(\nu+1)}\Big(\frac{x}{2}\Big)^\nu\, _{0}F_{1}\bigg[\,^{\displaystyle \,\,\,\,-}_{\displaystyle \nu+1}\,\,\bigg.\bigg|\,-{x^2}/{4}\,\bigg]

\end{equation} 

"--- функция Бесселя первого рода;

\begin{equation}

\label{monako:BesselI}

I_\nu (x)=\frac{1}{\Gamma(\nu+1)}\Big(\frac{x}{2}\Big)^\nu\, _{0}F_{1}\bigg[\,^{\displaystyle \,\,\,\,-}_{\displaystyle \nu+1}\,\,\bigg.\bigg|\,{x^2}/{4}\,\bigg]

\end{equation}

"--- модифицированная функция Бесселя первого рода;

\begin{equation}

\label{monako:GGamma}

\Gamma(\nu,x)=\int _x^\infty\xi^{\nu-1}\exp(-\xi)\,d\xi = \Gamma(\nu)-\frac{x^\nu}{\nu}\, _{1}F_{1}\Big[\,^{\displaystyle \,\,\,\,\nu}_{\displaystyle \nu+1}\,\,\Big.\Big|\,-x\,\Big]

\end{equation}

"--- неполная гамма-функция;

\begin{multline}

\label{monako:ierfc}

\mathop{\rm i^{\textit{n}}erfc}(x)=\frac{1}{2^n\,\Gamma(n / 2+1)}\, _{1}F_{1}\bigg[\,^{\displaystyle -n / 2}_{\displaystyle \,\,\,1 / 2}\,\,\bigg.\bigg|\,-x^2\,\bigg] - \\-\frac{2x}{2^{n}\,\Gamma((n+1) / 2)}\, _{1}F_{1}\bigg[\,^{\displaystyle (1-n) / 2}_{\displaystyle \,\,\,\,\,\,3 / 2}\,\,\bigg.\bigg|\,-x^2\,\bigg]

\end{multline}

"--- кратный дополнительный интеграл вероятности ($ n\in\mathbb{N} $).



Для гипергеометрической функции аргументов $x$, $y$:

$$F_{l:m;n}^{p:q;k}\bigg[\,^{\displaystyle (a_p):\,(b_q)\,;\,(c_k)}_{\displaystyle (d_l)\,:(f_m);(g_n)}\,\,\bigg.\bigg|\,x,\,y\bigg]=\sum_{r,s=0}^\infty \frac{\prod\limits_{\scriptscriptstyle j=1}^{\scriptscriptstyle p} (a_j)_{r+s}  \prod\limits_{\scriptscriptstyle j=1}^{\scriptscriptstyle q} (b_j)_{r}  \prod\limits_{\scriptscriptstyle j=1}^{\scriptscriptstyle k} (c_j)_{s}}{\prod\limits_{\scriptscriptstyle j=1}^{\scriptscriptstyle l} (d_j)_{r+s}  \prod\limits_{\scriptscriptstyle j=1}^{\scriptscriptstyle m} (f_j)_{r}  \prod\limits_{\scriptscriptstyle j=1}^{\scriptscriptstyle n} (g_j)_{s}}\frac{x^r y^s}{r!\,s!},$$

называемой функцией Кампе де Ферье (Kamp\'e de F\'eriet) \cite{manako:kampe}, можно записать частные случаи:

\begin{equation}

\label{monako:kampe1}

F_{1:1;0}^{1:1;0}\bigg[\,^{\displaystyle a:\,b\,\,;-}_{\displaystyle d:\,f\,;-}\,\,\bigg.\bigg|\,x,\,y\bigg]=\sum_{k,m=0}^\infty\frac{(a)_{k+m}(b)_k}{(d)_{k+m}(f)_k}\frac{x^k y^m}{k!\,m!},

\end{equation}

\begin{equation}

\label{monako:kampe2}

F_{1:1;1}^{2:0;0}\bigg[\,^{\displaystyle a_1,a_2:-\,;-}_{\displaystyle \,\,\,\,\,\, d\,\,\,\,\,:\,f\,;\,g}\,\,\bigg.\bigg|\,x,\,y\bigg]=\sum_{k,m=0}^\infty\frac{(a_1)_{k+m}(a_2)_{k+m}}{(d)_{k+m}(f)_k (g)_m}\frac{x^k y^m}{k!\,m!}.

\end{equation}

Выпишем функцию Гумберта  (Humbert)

\begin{equation}

\label{monako:Humbert}

\Psi_2(a\,;f,g\,;\,x,y)=\sum_{k,m=0}^\infty\frac{(a)_{k+m}}{(f)_k(g)_m}\frac{x^k y^m}{k!\,m!}=F_{0:1;1}^{1:0;0}\bigg[\,^{\displaystyle a:-\,;-}_{\displaystyle -\!:\,f\,;\,g}\,\,\bigg.\bigg|\,x,\,y\bigg]

\end{equation}

и функцию Аппеля (Appel)

\begin{multline}

\label{monako:Appel}

F_4(a_1,a_2;f,g\,;x,y)=\sum_{k,m=0}^\infty\frac{(a_1)_{k+m}(a_2)_{k+m}}{(f)_k(g)_m}\frac{x^k y^m}{k!\,m!}= \\ =F_{0:1;1}^{2:0;0}\bigg[\,^{\displaystyle a_1,a_2:-\,;-}_{\displaystyle \,\,\,\,\,-\,\,\,\,\,:\,f\,;g}\,\,\bigg.\bigg|\,x,\,y\bigg],

\end{multline}

которые также являются частными случаями функции  Кампе де Ферье.



Кроме перечисленных выше функций далее также будет использована общая гипергеометрическая функция Сриваставы (Srivastava), представимая в виде тройного гипергеометрического ряда \cite{manako:srivastava}:

\begin{multline*}

F^{(3)}\bigg[\,^{\displaystyle (a)::\,(b);(b');(b''):\,(d);(d');(d'')}_{\displaystyle (e)::(g);(g');(g''):(h);(h');(h'')}\,\,\bigg.\bigg|\,x,\,y,\,z\bigg]= \sum_{n,k,m=0}^\infty \frac{\prod\limits_{\scriptscriptstyle j=1}^{\scriptscriptstyle A} (a_j)_{n+k+m}  }{\prod\limits_{\scriptscriptstyle j=1}^{\scriptscriptstyle E} (e_j)_{n+k+m}  } \times\\

\times\frac{\prod\limits_{\scriptscriptstyle j=1}^{\scriptscriptstyle B} (b_j)_{n+k} \prod\limits_{\scriptscriptstyle j=1}^{\scriptscriptstyle B'} (b'_j)_{k+m} \prod\limits_{\scriptscriptstyle j=1}^{\scriptscriptstyle B''} (b''_j)_{n+m} \prod\limits_{\scriptscriptstyle j=1}^{\scriptscriptstyle D} (d_j)_{n} \prod\limits_{\scriptscriptstyle j=1}^{\scriptscriptstyle D'} (d'_j)_{k}  \prod\limits_{\scriptscriptstyle j=1}^{\scriptscriptstyle D''} (d''_j)_{m} }                                                                      {\prod\limits_{\scriptscriptstyle j=1}^{\scriptscriptstyle G} (g_j)_{n+k} \prod\limits_{\scriptscriptstyle j=1}^{\scriptscriptstyle G'} (g'_j)_{k+m} \prod\limits_{\scriptscriptstyle j=1}^{\scriptscriptstyle G''} (g''_j)_{n+m} \prod\limits_{\scriptscriptstyle j=1}^{\scriptscriptstyle H} (h_j)_{n} \prod\limits_{\scriptscriptstyle j=1}^{\scriptscriptstyle H'} (h'_j)_{k}  \prod\limits_{\scriptscriptstyle j=1}^{\scriptscriptstyle H''} (h''_j)_{m} } \frac{x^ny^kz^m}{n!\,k!\,m!}.

\end{multline*}

Здесь символ $ (a) $ обозначает множество чисел, состоящее из $ A $ элементов, символ $ (b) $ "--- множество из $ B $ элементов и~т.д. 

%Функция Сриваставы имеет абсолютную сходимость, если выполнена следующая система неравенств \cite{manako:srivastava}:

%$$

%\begin{array}{c}

%1+E+G+G''+H-A-B-B''-D\geq0,\\ 

%1+E+G+G'+H'-A-B-B'-D'\geq0,\\ 

%1+E+G'+G''+H''-A-B'-B''-D''\geq0.

%\end{array} 

%$$

Частным случаем функции Сриваставы является  функция 

\begin{multline}

\label{monako:srivastava}

F^{(3)}\bigg[\,^{\displaystyle a::\,\,b;-;-:\,-;\,\,\,-;-}_{\displaystyle e::-;-;-:\,\,\,h;\,\,h';-}\,\,\bigg.\bigg|\,x,\,y,\,z\bigg]=\\ =\sum_{n,k,m=0}^\infty\frac{(a)_{n+k+m}}{(e)_{n+k+m}}\frac{(b)_{n+k}}{(h)_{n}(h')_{k}}\frac{x^ny^kz^m}{n!\,k!\,m!}.

\end{multline}

%сходящаяся абсолютно при любых значениях аргументов  $x$,  $y$,  $z$, поскольку для неё 

%$ A=B=E=H=H'=1,$

%$ B'=B''=D=D'=D''=G=G'=G''=H''=0.$



\smallskip



\Section[n]{Расчёт нестационарной температуры} 

Метод мгновенных точечных источников \cite{manako:eger} решение краевой задачи  \eqref{monako:conductivity}, \eqref{monako:conditions} позволяет   записать в следующем виде:

\begin{multline}

T=2q\frac{a}{\lambda}\int _0^t\frac{dt'}{(4 \pi at')^{3/2}} \int _0^R r'\,dr'\times\\

\label{monako:triple}

\times\int _0^{2\pi}\exp\bigg(-\frac{(r\cos \varphi-r' \cos \varphi')^2+(r\sin \varphi-r' \sin \varphi')^2+z^2}{4at'}  \bigg)\,d\varphi'.

\end{multline}



\newpage



При $\nu=0$  для функции \eqref{monako:BesselI} имеет место \cite{manako:abramovits} интегральное представление 

$$I_0(x)=\frac{1}{2\pi}\int _0^{2\pi}\exp(x\cos \xi)\,d\xi,$$

позволяющее записать

%\begin{multline*}

%\int _0^{2\pi}\exp\bigg(-\frac{(r\cos \varphi-r' \cos \varphi')^2+(r\sin \varphi-r' \sin \varphi')^2+z^2}{4at'}  \bigg)\,d\varphi'=\\

%=\int _0^{2\pi}\exp\bigg(-\frac{(r^2+r'^2-2rr'\cos(\varphi'-\varphi)+z^2}{4at'}  \bigg)\,d\varphi'= \\ =2\pi\exp\bigg(-\frac{r^2+r'^2+z^2}{4at'}\bigg)\,I_0\bigg(\frac{rr'}{2at'}\bigg).

%\end{multline*}

%Тогда с учетом \eqref{monako:BesselI} представим выражение 

\eqref{monako:triple} в виде 

\begin{multline}

\label{monako:double}

T=4\pi q\frac{a}{\lambda}\int _0^t\frac{dt'}{(4 \pi at')^{3/2}} \times \\ \times  \int _0^R\exp\bigg(-\frac{r^2+r'^2+z^2}{4at'}\bigg)\,_{0}F_{1}\Bigg[\,^{\displaystyle -}_{\displaystyle \,1}\,\,\bigg.\bigg|\,\bigg(\frac{rr'}{4at'}\bigg)^2\,\Bigg] r'\,dr'.

\end{multline}

Представляя экспоненту и обобщённую гипергеометрическую функцию $ _0F_1 $ в виде рядов, можно вычислить следующий определённый интеграл:

$$

%\begin{multline*}

\int _0^b\exp(-a\xi)\,_{0}F_{1}\Big[\,^{\displaystyle -}_{\displaystyle \,1}\,\,\Big.\Big|\,\xi\,\Big]\,d\xi=

%\sum_{n,k=0}^\infty\frac{(-1)^n}{n!}\frac{a^n}{(1)_k k!}\Bigg[\int _0^b\xi^{n+k}\,d\xi\Bigg]=\\

%=\sum_{n,k=0}^\infty\frac{(-1)^n}{n!}\frac{a^n}{(1)_k k!}\frac{b^{n+k+1}}{(n+k+1)}=b\sum_{m=0}^\infty\sum_{k=0}^m\frac{(-1)^{m+k}}{(m-k)!}\frac{a^{m-k}}{(1)_k k!}\frac{b^{m}}{(m+1)}=\\

%= b\sum_{m=0}^\infty\frac{(-1)^{m}}{m!}\frac{(ab)^m}{(m+1)}\,_{1}F_{1}\bigg[\,^{\displaystyle -m}_{\displaystyle \,\,\,\,\,1}\bigg.\bigg|\,\frac{1}{a}\,\bigg]= 

%\\=b\exp\bigg(\frac{1}{a}\bigg)\sum_{m=0}^\infty\frac{(-1)^{m}}{m!}\frac{(ab)^m}{(m+1)}\,_{1}F_{1}\bigg[\,^{\displaystyle 1+m}_{\displaystyle \quad 1}\bigg.\bigg|\,\frac{-1}{a}\,\bigg]=\\

%=b\exp\bigg(\frac{1}{a}\bigg)\sum_{m,k=0}^\infty\frac{(1+m)_{k}}{(1)_k k!\,m!}\frac{(-a)^{-k}(-ab)^m}{(m+1)}=\\ =b\exp\bigg(\frac{1}{a}\bigg)\sum_{m,k=0}^\infty\frac{(1)_{m+k}}{(1)_k(2)_m }\frac{(-a)^{-k}(-ab)^m}{k!\,m!}=

%

b\exp\big({1}/{a}\big) \Psi_2\big(1;1,\,2;{-1}/{a},-ab\big).

%\end{multline*}

$$

При вычислении интеграла учитывались выражения \eqref{monako:hypergeom}, \eqref{monako:Humbert}, а также первая теорема Куммера \cite{manako:abramovits,manako:prud3}.

%:

%$$_{1}F_{1}\Big[\,^{\displaystyle a}_{\displaystyle b}\,\,\Big.\Big|\,x\,\Big]=\exp\big(x\big)\,_{1}F_{1}\Big[\,^{\displaystyle b-a}_{\displaystyle \quad  b}\,\,\Big.\Big|\,-x\,\Big].$$

Используя полученное равенство,

%$$\int _0^b\exp(-a\xi)\,_{0}F_{1}\Big[\,^{\displaystyle -}_{\displaystyle \,1}\,\,\Big.\Big|\,\xi\,\Big]\,d\xi=b\exp\bigg(\frac{1}{a}\bigg)\,\Psi_2\bigg(1;1,\,2;\frac{-1}{a},-ab\bigg),$$

представим \eqref{monako:double} в виде

$$T(r, z, t)=2\pi r^2 q\frac{a}{\lambda}\int _0^t \exp\bigg(-\frac{z^2}{4at'}\bigg)\,\Psi_2\bigg(1;1,\,2;-\frac{r^2}{4at'},-\frac{R^2}{4at'}\bigg)  \frac{dt'}{(4 \pi at')^{3/2}}.$$

Отсюда, вводя  безразмерные величины

\begin{equation}

\label{monako:less}

\theta={\lambda T}/({qR}),\quad \tau={4at}/{R^2},\quad \sigma={z}/{R},\quad \rho={r}/{R},

\end{equation}

имеем

\begin{equation}

\label{monako:theta}

\theta(\rho, \sigma, \tau)=\frac{1}{2\sqrt{\pi}}\int _{1/\tau}^\infty \xi^{\,-1/2}\exp\big(-\sigma^2\xi\big)\,\Psi_2\big(1;1,\,2;-\rho^2\xi,-\xi\big)\,d\xi.

\end{equation}

Раскладывая в ряд функцию \eqref{monako:Humbert} и учитывая \eqref{monako:hypergeom}, \eqref{monako:GGamma}, \eqref{monako:Appel}, \eqref{monako:srivastava}, вычислим интеграл в  \eqref{monako:theta} и получим

%$$ \theta=\frac{1}{2\sqrt{\pi}}\int _{1/\tau}^\infty \xi^{\,-1/2}\exp\big(-\sigma^2\xi\big)\,\Psi_2\big(1;1,\,2;-\rho^2\xi,-\xi\big)\,d\xi= $$

%$$=\frac{1}{2\sqrt{\pi}}\sum_{n,k=0}^\infty\frac{(1)_{n+k}}{(1)_{n}(2)_{k}}\frac{(-1)^{n+k}\rho^{2n}}{n!\,k!}\int _{1/\tau}^\infty \xi^{n+k-1/2}\exp\big(-\sigma^2\xi\big)\,d\xi =$$

%$$=\frac{1}{2\sigma\sqrt{\pi}}\sum_{n,k=0}^\infty\frac{(1)_{n+k}}{(1)_{n}(2)_{k}n!\,k!}\bigg(\frac{-1}{\sigma ^2}\bigg)^{n+k}\rho^{2n}\Gamma \bigg(1 / 2+n+k,\frac{\sigma^2}{\tau}\bigg)= $$

%$$ =\frac{1}{2\sigma}\sum_{n,k=0}^\infty\frac{(1 / 2)_{n+k}(1)_{n+k}}{(1)_{n}(2)_{k}n!\,k!}\bigg(\frac{-1}{\sigma ^2}\bigg)^{n+k}\rho^{2n}- $$

%$$- \frac{1}{\sqrt{\tau \pi}}\sum_{n,k=0}^\infty \frac{(1)_{n+k}(1 / 2)_{n+k}\,\,\,\rho^{2n}}{(1)_{n}(2)_{k}(3 / 2)_{n+k}n!\,k!}\bigg(\frac{-1}{\tau}\bigg)^{n+k}  \,_{1}F_{1}\bigg[\,^{\displaystyle 1/2+n+k}_{\displaystyle 3/2+n+k}\,\,\bigg.\bigg|\,\frac{-\sigma^2}{\tau}\,\bigg]=$$

%$$=\frac{1}{2\sigma}F_4\bigg(1 / 2,\,1;\,1,\,2;\frac{-\rho^2}{\sigma^2},\frac{-1}{\sigma^2}\bigg)- $$

%$$- \frac{1}{\sqrt{\tau \pi}}\sum_{n,k,m=0}^\infty \frac{(1 / 2)_{n+k+m}(1)_{n+k}}{(3 / 2)_{n+k+m}(1)_{n}(2)_{k}n!\,k!\,m!}  \bigg(\frac{-\rho^2}{\tau}\bigg)^{n} \bigg(\frac{-1}{\tau}\bigg)^{k} \bigg(\frac{-\sigma^2}{\tau}\bigg)^{m}  =$$

%$$=\frac{1}{2\sigma}F_4\bigg(1 / 2,\,1;\,1,\,2;\frac{-\rho^2}{\sigma^2},\frac{-1}{\sigma^2}\bigg)- \frac{1}{\sqrt{\tau \pi}} F^{(3)}\bigg[\,^{\displaystyle 1 / 2::\,1\,;-;-:\,-;-;-}_{\displaystyle 3 / 2::-;-;-:\,\,1\,;\,2\,;-}\,\,\bigg.\bigg|\,\frac{-\rho^2}{\tau},\,\frac{-1}{\tau},\,\frac{-\sigma^2}{\tau}\bigg].$$

%

%Таким образом, формула 

\begin{multline}

\theta(\rho,\sigma,\tau)=\frac{1}{2\sigma}F_4\big(1 / 2,\,1;\,1,\,2;{-\rho^2}/{\sigma^2}, {-1}/{\sigma^2}\big)-\\

\label{monako:theta1}

- \frac{1}{\sqrt{\tau \pi}} F^{(3)}\bigg[\,^{\displaystyle 1 / 2::\,1\,;-;-:\,-;-;-}_{\displaystyle 3 / 2::-;-;-:\,\,1\,;\,2\,;-}\,\,\bigg.\bigg|\,{-\rho^2}/{\tau},\,{-1}/{\tau},\,{-\sigma^2}/{\tau}\bigg].

\end{multline}





Функция \eqref{monako:theta1} является аналитическим решением следующей безразмерной краевой задачи:

$$

\frac{\partial^2\theta}{\partial \rho^2}+\frac{1}{\rho}\frac{\partial \theta}{\partial \rho}+\frac{\partial^2\theta}{\partial \sigma^2}=4\frac{\partial \theta}{\partial \tau},

$$

$$

\frac{\partial \theta}{\partial \sigma}\Big|_{\sigma\to  0}=-H(1-\rho), \quad 

\frac{\partial \theta}{\partial \rho}\Big|_{\rho \to  0}= \theta\big|_{\rho\to \infty}=\theta\big|_{\sigma\to \infty}= \theta\big|_{\tau\to 0}=0.

$$

%

%

%

%есть представление через функцию Аппеля \eqref{monako:Appel} и общую гипергеометрическую функцию Сриваставы \eqref{monako:srivastava} для нестационарной температуры полубесконечного тела, нагреваемого неподвижным круговым поверхностным источником тепла. Из \eqref{monako:conductivity}, \eqref{monako:conditions}, \eqref{monako:less} следует, что функция \eqref{monako:theta1} является аналитическим решением в полуограниченном пространстве следующей краевой задачи (все параметры безразмерны): 

%$$\frac{\partial^2\theta}{\partial \rho^2}+\frac{1}{\rho}\frac{\partial \theta}{\partial \rho}+\frac{\partial^2\theta}{\partial \sigma^2}=4\frac{\partial \theta}{\partial \tau},$$

%$$\bigg.\frac{\partial \theta}{\partial \sigma}\bigg|_{\sigma\to  0}=-\,H(1-\rho), \qquad \qquad \bigg.\frac{\partial \theta}{\partial \rho}\bigg|_{\rho \to  0}=\bigg.\theta\bigg|_{\rho\to \infty}=\bigg.\theta\bigg|_{\sigma\to \infty}=\bigg.\theta\bigg|_{\tau\to 0}=0.$$

Заметим, что в работе \cite{manako:eger} решение этой краевой задачи, полученное с помощью интегральных преобразований Ханкеля и Фурье, представлено в~виде 

\begin{multline*}

\theta(\rho,\sigma,\tau)=\frac{1}{2}\int _0^\infty J_0(\rho\xi)\,J_1(\xi)\bigg[\exp(-\sigma\xi)\mathop{\rm erfc}\Big(\frac{\sigma}{\sqrt{\tau}}-\frac{\xi}{2}\sqrt{\tau}\Big) -\\-\exp(\sigma\xi)\mathop{\rm erfc}\Big(\frac{\sigma}{\sqrt{\tau}}+\frac{\xi}{2}\sqrt{\tau}\Big)\bigg]\frac{d\xi}{\xi}.

\end{multline*}

Очевидно, что это решение совпадает с  \eqref{monako:theta1}.



\smallskip



\Section[n]{Частные случаи формулы (\ref{monako:theta1})} Сделаем ряд замечаний. 



$1^\circ$. Первое слагаемое в правой части \eqref{monako:theta1} задаёт стационарное температурное поле полубесконечного тела, нагретого круговым источником:

\begin{equation}

\label{monako:16}

\theta (\rho,\sigma,\infty)=\frac{1}{2\sigma}F_4\big(1 / 2,\,1;\,1,\,2;{-\rho^2}/{\sigma^2},{-1}/{\sigma^2}\big).

\end{equation}

Отметим, что в работе \cite{manako:TVT} найдено стационарное температурное поле для гауссова источника тепла:

\begin{equation}

\label{monako:300}

\theta_{\text {гаус}}(\rho,\sigma,\infty)=\Psi_2\big(1 / 2;1 / 2,\,1;\sigma^2,-\rho^2\big)-\frac{2\sigma}{\sqrt{\pi}}\Psi_2\big(1;3 / 2,\,1;\sigma^2,-\rho^2\big).

\end{equation}



\begin{wraptable}[12]{O}{.5\textwidth}

\small \centering \vspace{-2mm}

\begin{tabular}{c|c|c|c}

\hline 

\footnotesize $ \sigma $ & 

\footnotesize $ \rho $ & 

\footnotesize $ \theta (\rho,\sigma,\infty)$ & 

\footnotesize $ \theta_{\text {гаус}}(\rho,\sigma,\infty)$ \\  \hline

\rule{0mm}{12pt}%

0,01& 0,01 & 0,990 & 0,989  \\ 

0,08& 0,1 & 0,921 & 0,912 \\ 

0,1& 0,01 & 0,905 & 0,896 \\ 

0,08& 0,6 & 0,827 & 0,783 \\ 

0,4& 0,3 & 0,659 & 0,652 \\ 

0,5& 0,9 & 0,473 & 0,496 \\ 

1,0& 0,0 & 0,414 & 0,428 \\ 

1,0& 0,6 & 0,384 & 0,401 \\ 

1,6& 0,6 & 0,274 & 0,294 \\ 

1,6& 1,2 & 0,242 & 0,265 \\ \hline

\end{tabular}



\end{wraptable}



\noindent 

Расчёт по последним двум формулам даёт ожидаемый результат: модели кругового и гауссова источников тепла предсказывают близкие по значениям температурные поля (см. таблицу).







Заметим, что в соответствии с~\cite{manako:srivastava2} выражение \eqref{monako:300} можно представить в компактной форме:

$$ \theta_{\text {гаус}}(\rho,\sigma,\infty)=\sqrt{{2}({\rho\pi})^{-1}}V_{0,0}(\rho,\sigma), $$

где $V_{\mu,\nu}(x,y)$ "--- обобщённая функция Фойгта (Voigt). %равна (при условии $ {\rm Re}(\mu+\nu)>-1 $)

%$$V_{\mu,\nu}(x,y)= \sqrt{\frac{x}{2}}\int _0^\infty\xi^{\mu}\exp\Big(-y\xi-\frac{\xi^2}{4}\Big)\,J_{\nu}(x\xi)\,d\xi=$$

%$$=\frac{2^{\mu-1/2}\,x^{\nu+1/2}}{\Gamma(\nu+1)}\bigg\{\Gamma\Big(\frac{\mu+\nu+1}{2}\Big)\,\Psi_2\Big(\frac{\mu+\nu+1}{2};\,\frac{1}{2},\,\nu+1;\,y^2,\,-x^2\Big)-$$

%$$-2y \,\Gamma\Big(\frac{\mu+\nu+2}{2}\Big)\,\Psi_2\Big(\frac{\mu+\nu+2}{2};\,\frac{3}{2},\,\nu+1;\,y^2,\,-x^2\Big)\bigg\}.$$



\smallskip



$2^\circ$. Полагая в формуле \eqref{monako:theta1} $ \rho=0 $ и учитывая уменьшение количества аргументов для функции Аппеля (см. \eqref{monako:hypergeom} и \eqref{monako:Appel}):

$$ F_4\big(1 / 2,\,1;\,1,\,2;0,\,y\big)=  _{2}F_{1}\bigg[\,^{\displaystyle 1 / 2,\,1}_{\displaystyle \quad \,\, 2}\,\,\bigg.\bigg|\,y\,\bigg] $$

и для функции Сриваставы (см. \eqref{monako:kampe1} и \eqref{monako:srivastava}):

$$F^{(3)}\bigg[\,^{\displaystyle 1 / 2::\,1\,;-;-:\,-;-;-}_{\displaystyle 3 / 2::-;-;-:\,\,1\,;\,2\,;-}\,\,\bigg.\bigg|\,0,\,y,\,z\bigg]=F_{1:1;0}^{1:1;0}\bigg[\,^{\displaystyle 1 / 2:1;-}_{\displaystyle 3 / 2:2;-}\,\,\bigg.\bigg|\,y,\,z\bigg],$$

нестационарное температурное поле по оси луча лазера можно записать как

\begin{equation}

\label{monako:18} \!

\theta(0, \sigma,\tau)=\frac{1}{2\sigma}{}_{2}F_{1}\bigg[{}^{\displaystyle 1 / 2,\,1}_{\displaystyle \quad \,\, 2} \bigg|\,{-1}/{\sigma^2}\,\bigg]-\frac{1}{\sqrt{\tau\pi}}F_{1:1;0}^{1:1;0}\bigg[\,^{\displaystyle 1 / 2:1;-}_{\displaystyle 3 / 2:2;-}\,\,\bigg.\bigg|\,{-1}/{\tau}, {-\sigma^2}/{\tau}\bigg]. \!

\end{equation}

%Численные расчеты при различных значениях $ \sigma $ и $ \tau $ показывают, что результат \eqref{monako:18} совпадает с  выражением \cite{manako:eger}

Формула \eqref{monako:18} совпадает с  формулой \cite{manako:eger}

\begin{equation}

\label{monako:19}

\theta(0,\sigma,\tau)=\sqrt{\tau}\mathop{\rm ierfc}\big({\sigma}/{\sqrt{\tau}}\big)-\sqrt{\tau}\mathop{\rm ierfc}\sqrt{({\sigma^2+1})/{\tau}}.

\end{equation}

Однако формула \eqref{monako:18} имеет по сравнению с \eqref{monako:19} преимущество: в неё входит максимальная (стационарная) температура нагрева на оси лазерного луча в явном виде:

$$ \theta (0,\sigma,\infty)=\frac{1}{2\sigma}\,_{2}F_{1}\bigg[\,^{\displaystyle 1 / 2,\,1}_{\displaystyle \quad \,\, 2}\,\,\bigg.\bigg|\,{-1}/{\sigma^2}\,\bigg]=\sqrt{\sigma^2+1}-\sigma. $$ 



\smallskip



$3^\circ$. Рассмотрим нестационарное температурное поле на поверхности полубесконечного тела \cite{manako:eger}: 

\begin{equation}

\label{monako:200}

\theta(\rho,0,\tau)=\int _0^\infty\,J_0(\rho\xi)\,J_1(\xi)\frac{d\xi}{\xi}-\int _0^\infty\,J_0(\rho\xi)\,J_1(\xi)\mathop{\rm erfc}\bigg(\frac{\xi}{2}\sqrt{\tau}\bigg)\frac{d\xi}{\xi}.

\end{equation}

Стационарную часть \eqref{monako:200} можно записать в виде 

\begin{multline}

\label{monako:201}

\theta (\rho, 0, \infty)=\int _0^\infty J_0(\rho\xi) J_1(\xi)\frac{d\xi}{\xi}=  \\ =

\left\{

\begin{array}{cl}  

{}_{2}F_{1}\Big[\,^{ -1 / 2,\,1 / 2}_{\qquad\,\,\,  1}\,\,\big|\,\rho^2\,\Big],&\text{при $ \rho\leq1 $};\\

{({2\rho})^{-1}} {}_{2}F_{1}\Big[\,^{1 / 2,\,1 / 2}_{\quad \,\,\,\, 2}\,\,\big|\,\rho^{-2}\,\Big],&\text{при $ \rho\geq1 $.}\end{array}\right.

\end{multline}

Значение последнего интеграла, входящего в правую часть \eqref{monako:200}, %неизвестно \cite{manako:prud3,manako:prud2}. Однако, его 

можно записать в аналитическом виде, если в формуле \eqref{monako:theta1} для функции Сриваставы $ F^{(3)} $ положить $ \sigma=0 $. 

При этом функция $ F^{(3)} $ за счёт уменьшения количества аргументов вырождается в функцию Кампе де Ферье \eqref{monako:kampe2}:

$$ \int _0^\infty J_0(\rho\xi) J_1(\xi)\mathop{\rm erfc}\bigg(\frac{\xi}{2}\sqrt{\tau}\bigg)\frac{d\xi}{\xi}=\frac{1}{\sqrt{\tau\pi}}F_{1:1;1}^{2:0;0}\bigg[\,^{\displaystyle 1 / 2,\,1:-;-}_{\displaystyle \,\,\,\,3 / 2\,\,\,:\,1\,;\,2}\,\,\bigg.\bigg| {-\rho^2}/{\tau}, {-1}/{\tau}\bigg].$$

Следовательно, формула \eqref{monako:200} представима в виде

\begin{equation}

\label{monako:202}

\theta(\rho, 0, \tau)=\theta (\rho,0,\infty)-\frac{1}{\sqrt{\tau\pi}}F_{1:1;1}^{2:0;0}\bigg[\,^{\displaystyle 1 / 2,\,1:-;-}_{\displaystyle \,\,\,\,3 / 2\,\,\,:\,1\,;\,2}\,\,\bigg.\bigg|\,{-\rho^2}/{\tau}, {-1}/{\tau}\bigg],

\end{equation}

где $ \theta (\rho,0,\infty) $ определяется из \eqref{monako:201}.



$4^\circ$. Рассмотрим нестационарное температурное поле на цилиндрической поверхности $ \rho=1 $, ориентированной соосно лазерному лучу:

\begin{multline*}

 \theta (1,\sigma,\tau)=\frac{1}{2\sigma}F_4\big(1 / 2,\,1;\,1,\,2; {-1}/{\sigma^2},{-1}/{\sigma^2}\big)- \\

- \frac{1}{\sqrt{\tau \pi}} F^{(3)}\bigg[\,^{\displaystyle 1 / 2::\,1\,;-;-:\,-;-;-}_{\displaystyle 3 / 2::-;-;-:\,\,1\,;\,2\,;-}\,\,\bigg.\bigg| {-1}/{\tau}, {-1}/{\tau}, {-\sigma^2/}{\tau}\bigg].

\end{multline*}

Данное выражение можно привести к виду

\begin{multline}

\theta (1,\sigma,\tau)=\frac{1}{2\sigma}\,_{4}F_{3}\bigg[\,^{\displaystyle 1 / 2,\,1,\,1,\,3 / 2}_{\displaystyle \quad \quad \,\,1,\,2,\,2}\,\,\bigg.\bigg| {-4}/{\sigma^2}\,\bigg]-\\

\label{monako:20}

- \sqrt{\frac{\tau}{4\pi}}\Bigg\{F_{1:1;0}^{1:1;0}\bigg[\,^{\displaystyle -1 / 2:1 / 2;-}_{\displaystyle \,\,\, 1 / 2 \,\,:\,\,\,\,1\,\,;-}\,\,\bigg.\bigg| {-4}/{\tau}, {-\sigma^2}/{\tau}\bigg]- {}_{1}F_{1}\bigg[\,^{\displaystyle -1 / 2}_{\displaystyle \,\,\,1 / 2}\,\,\bigg|\, {-\sigma^2}/ {\tau}\,\bigg]   \Bigg\},

\end{multline}

из  которого легко выражается стационарная часть:

\begin{multline*}

 \theta (1,\sigma,\infty)=\frac{1}{2\sigma}\,_{4}F_{3}\bigg[\,^{\displaystyle 1 / 2,\,1,\,1,\,3 / 2}_{\displaystyle \quad \quad \,\,1,\,2,\,2}\,\,\bigg.\bigg|\,{-4}/{\sigma^2}\,\bigg]=\\ =\sqrt{\frac{\sigma^2+4}{4}}\,_{2}F_{1}\bigg[\,^{\displaystyle -1 / 2,\,1 / 2}_{\displaystyle \quad \quad 1}\,\,\bigg.\bigg|\,{4}/({\sigma^2+4}) \bigg]-\frac{\sigma}{2}.

 \end{multline*}

Полученный результат совпадает с результатом работы \cite{manako:pinsker}, в которой гипергеометрическая функция Гаусса $ _{2}F_{1} $ выражается через полный эллиптический интеграл второго рода.





Если в выражении \eqref{monako:20} положить $ \sigma=0 $, то за счёт уменьшения количества аргументов в функции Кампе де Ферье  можно записать формулу

\begin{equation}

\label{monako:21}

\theta (1,0,\tau)=\frac{2}{\pi}-\sqrt{\frac{\tau}{4\pi}}\Bigg\{\,_{1}F_{1}\bigg[\,^{\displaystyle -1 / 2}_{\displaystyle \quad \, 1}\,\,\bigg.\bigg| {-4}/{\tau}\,\bigg]-1\Bigg\},

\end{equation}

которая позволяет определить температуру на краю поверхностного пятна нагрева. 

Температура в центре этого пятна вычисляется по формуле

$$

\theta (0,0,\tau)=1-\frac{1}{\sqrt{\tau\pi}}\,_{2}F_{2}\bigg[\,^{\displaystyle 1 / 2,\,1}_{\displaystyle 3 / 2,\,2}\,\,\bigg.\bigg|\, {-1}/{\tau}\,\bigg].

$$

Если здесь применить формулу понижения порядка обобщённой гипергеометрической функции \cite{manako:prud3}

$$ _{p}F_{q}\bigg[\,^{\displaystyle (a_{p-1}),\,1}_{\displaystyle \,(b_{q-1}),\,2}\,\,\bigg.\bigg|\,x\,\bigg]= x^{-1}\frac{\prod\limits_{\scriptscriptstyle j=1}^{\scriptscriptstyle q-1} (b_j-1)}{\prod\limits_{\scriptscriptstyle k=1}^{\scriptscriptstyle p-1} (a_k-1)}\Bigg\{\,_{p-1}F_{q-1}\bigg[\,^{\displaystyle (a_{p-1})-1}_{\displaystyle \,(b_{q-1})-1}\,\,\bigg.\bigg|\,x\,\bigg] -1\Bigg\},$$

то последнему выражению можно придать вид

\begin{equation}

\theta (0,0,\tau)=1-\sqrt{\frac{\tau}{\pi}}\Bigg\{\,_{1}F_{1}\bigg[\,^{\displaystyle -1 / 2}_{\displaystyle \,\,\,\, 1 / 2}\,\,\bigg.\bigg|\,{-1}/{\tau}\,\bigg]-1\Bigg\}.

\label{monako:new}

\end{equation}



Формулы \eqref{monako:21}, \eqref{monako:new} могут быть записаны в виде 

%

%Отсюда и \eqref{monako:21} с учетом известных \cite{manako:abramovits,manako:prud3} свойств гипергеометрической функции Куммера $ _1F_1 $ найдем 

\begin{gather}

\label{monako:101}

\theta (1, 0, \tau)=\frac{2}{\pi}-\sqrt{\frac{\tau}{4\pi}}\bigg\{\exp\bigg(\frac{-2}{\tau}\bigg)\bigg[\bigg(1+\frac{4}{\tau}\bigg)\,I_0\bigg(\frac{2}{\tau}\bigg)+\frac{4}{\tau}\,I_1\bigg(\frac{2}{\tau}\bigg) \bigg]-1\bigg\},

\\

\label{monako:100}

\theta (0, 0, \tau)=\mathop{\rm erfc}\bigg(\frac{1}{\sqrt{\tau}}\bigg)+\sqrt{\frac{\tau}{\pi}}\bigg[1-\exp\bigg(\frac{-1}{\tau}\bigg)\bigg],

\end{gather}

в котором они получены в \cite{manako:pinsker}. %Переход к размерным переменным в \eqref{monako:theta1} осуществляется с помощью соотношений \eqref{monako:less}.



%

%Таким образом, ранее установленные \cite{manako:pinsker} выражения \eqref{monako:100} и \eqref{monako:101} являются частными случаями формул \eqref{monako:18}, \eqref{monako:20}, вытекающих из наиболее общей формулы \eqref{monako:theta1}.  



\newpage



\Section[n]{Заключение} В работе получено аналитическое выражение \eqref{monako:theta1} для нестационарного температурного поля при нагреве полубесконечного тела лучом лазера, создающим на поверхности круговой неподвижный источник тепла. 

Первое слагаемое в \eqref{monako:theta1} описывает стационарное поле температур и выражается через функцию Аппеля \eqref{monako:Appel}, т.е. гипергеометрическую функцию двух аргументов. 

Второе слагаемое, учитывающее нестационарность процесса, представлено через гипергеометрическую функцию трёх аргументов (функция Сриваставы \eqref{monako:srivastava}). 

Функции \eqref{monako:Appel} и \eqref{monako:srivastava} допускают уменьшение количества аргументов в~трёх случаях: 

на оси лазерного луча (см. формулу \eqref{monako:18}), на свободной поверхности (см. формулу \eqref{monako:202}) и на круговой цилиндрической поверхности, являющейся продолжением границы зоны нагрева (см. формулу~\eqref{monako:20}). 



Рассмотрены стационарные температурные поля, образуемые круговым и гауссовым источниками. Существенных количественных различий в этих полях не обнаружено. 



\thanks{Автор выражает искреннюю благодарность В.~А.~Пинскеру, сделавшему ценные и полезные замечания после прочтения рукописи данной работы.}



\begin{thebibliography}{99}



\RBibitem{manako:pinsker}

\by Пинскер~В.\,А.

\paper Нестационарное температурное поле в полуограниченном теле, нагреваемом круговым поверхностным источником тепла

\jour Теплофизика высоких температур

\yr 2006

\vol 44

\issue 1

\pages 127--135

\transl англ. пер.:~

\paper Unsteady-state temperature field in a semi-infinite body heated by a disk surface heat source

\by Pinsker~V.\,A.

\jour High Temperature

\yr 2006

\vol 44

\issue 1

\pages 129--138



\Bibitem{manako:kampe}

\by  Appel~P.,  Kamp\'e de F\'eriet~J.

\book Fonctions Hyperg\'eom\'etriques et Hypersph\'eriques. Polyn\^omes d’Hermite

\publaddr Paris

\publ Gauthier-Villars

\yr 1926

\totalpages 202



\Bibitem{manako:srivastava}

\by  Srivastava~H.\,M.,  Manocha~H.\,L.

\book A  treatise  on  generating  functions 

\publaddr New York

\publ Halsted Press [John Wiley \& Sons, Inc.]

\yr 1984

\totalpages 569





\Bibitem{manako:eger}

\by Carslaw~H.\,S., Jaeger~J.\,C.

\book Conduction of Heat in Solids

\publaddr Oxford

\publ Clarendon Press

\yr 1959

\totalpages 510

\rtransl русск. пер.:~

\by Карслоу Г., Егер Д.

\book Теплопроводность твёрдых тел 

\publaddr М.

\publ Наука

\yr 1964

\totalpages 488



\Bibitem{manako:abramovits}

\book Handbook of Mathematical Functions with Formulas, Graphs, and Mathematical Tables

\yr 1972

\eds M.~Abramowitz, I.~A.~Stegun

\publ Dover

\publaddr New York

\totalpages 824

\rtransl

русск. пер.:~

\book Справочник по специальным функциям

\yr 1979

\eds М.~Абрамовиц, И.~Стиган

\publ Наука

\publaddr М.

\totalpages 832

\mathscinet{http://www.ams.org/mathscinet-getitem?mr=563328}



\RBibitem{manako:prud3}

\by Прудников А.\,П., Брычков Ю.\,А., Маричев О.\,И. 

\book Интегралы и ряды

\bookvol 3

\volinfo Специальные функции. Дополнительные главы 

\publaddr М.

\publ Физматлит

\yr 2003

\totalpages 688

\misctransl

\by Prudnikov~A.\,P. Brychkov~Yu.\,A., Marichev~O.\,I.

\book Integrals and series

\bookvol 3

\volinfo Special functions. Supplementary chapters

\publ Fizmatlit

\publaddr Moscow

\yr 2003

\totalpages 688



\RBibitem{manako:TVT}

\by Манако~В.\,В., Путилин~В.\,А., Камашев~А.\,В.

\paper Расчет температуры при нагреве неподвижным лазерным лучом

\jour Теплофизика высоких температур

\yr 2011

\vol 49

\issue 1

\pages 126--132

\transl

англ. пер.:~

\paper Calculation of temperature under heating by an immobile laser beam

\by Manako~V.\,V., Putilin~V.\,A., Kamashev~A.\,V.

\jour High Temperature

\yr 2011

\vol 49

\issue 1

\pages 127--134





\Bibitem{manako:srivastava2}

\by Srivastava~H.\,M., Miller~E.\,A.  

\paper A unified presentation of the Voigt functions

\jour Astrophys. and Space Sci.

\yr 1987

\vol 135

\issue 1

\pages 111--118





\end{thebibliography}





\noindent

\hskip6.5cm \thedate \par\noindent

\hskip6.5cm\therevdate



\pagebreak





\chead[\fancyplain{}{\scriptsize \sl M\,a\,n\,a\,k\,o~V.\,V.}]

{\fancyplain{}{\scriptsize \sl  A representation in terms of hypergeometric functions \dots}}





\makeengtitlem



\newpage



%\end{document}

